\section{Search-o1: recuperación autónoma en pleno proceso de razonamiento}

\subsection{Por qué los razonamientos largos son más propensos a tener vacíos de conocimiento}

\videofigure{figures/search-o1-intro.jpg}{Search-o1 agrega capacidad de recuperación autónoma a un large reasoning model, de modo que pueda completar conocimiento según lo necesite durante el razonamiento.}{00:46:56--00:48:02}

\videofigure{figures/reasoning-knowledge-gaps.jpg}{Cuanto más larga es la cadena de razonamiento, más oportunidades hay de encontrar entidades desconocidas, fórmulas, hechos experimentales o conocimiento interdisciplinario, y los vacíos tempranos siguen propagándose.}{00:48:02--00:49:28}

Un large reasoning model puede tener una capacidad deductiva sólida y aun así carecer del hecho necesario para completar cierto paso. Si el modelo oculta el vacío con expresiones ambiguas como «perhaps», «might» o «likely», el razonamiento posterior seguirá desarrollándose sobre una premisa errónea.

\begin{importantbox}{Un fallo de reasoning no siempre se debe a una capacidad de reasoning insuficiente}
Algunos fallos provienen de la falta de conocimiento; otros, de la incapacidad de razonar. El valor de Search-o1 está en distinguir ambos casos: ante un vacío de conocimiento, primero recupera información y, una vez obtenida la evidencia, evalúa si el modelo aún es incapaz de completar la deducción.
\end{importantbox}

\subsection{Por qué un RAG de una sola pasada no es suficiente}

\videofigure{figures/standard-rag-limitations.jpg}{Un RAG estándar normalmente solo recupera una vez al inicio, con la pregunta original, y no puede responder a las nuevas preguntas que surgen a mitad del razonamiento.}{00:49:28--00:50:24}

La pregunta original suele no poder expresar directamente la necesidad real de recuperación. Por ejemplo, un problema de química de varios pasos requiere primero identificar un producto intermedio y después consultar cierta propiedad de reacción; la consulta del segundo paso solo se conoce una vez completado el razonamiento del primer paso.

\begin{warningbox}{Más documentos pueden empeorar el reasoning}
Meter todos los documentos del top-$k$ en el context introduce redundancia, contradicciones y una dilución de la atención por contexto largo. El sistema de recuperación debe resolver a la vez «cuándo buscar», «qué buscar» y «qué conservar».
\end{warningbox}

\subsection{Agentic RAG y Reason-in-Documents}

\videofigure{figures/search-o1-framework.jpg}{Search-o1 añade búsqueda de múltiples rondas y razonamiento dentro de los documentos, sobre la base del vanilla reasoning y del RAG de una sola pasada.}{00:50:24--00:54:42}

La ronda $t$ se puede representar como:

$$
u_t=g(r_{\le t}),\quad
q_t=h(u_t),\quad
D_t=\operatorname{Retrieve}(q_t),\quad
z_t=\operatorname{Compress}(D_t,q_t),
$$

$$
r_{t+1}\sim p_\theta(\cdot\mid x,r_{\le t},z_t).
$$

\begin{itemize}
    \item $r_{\le t}$: el prefijo de reasoning actual;
    \item $u_t$: el punto de incertidumbre o vacío de conocimiento que identifica el modelo;
    \item $q_t$: la consulta de búsqueda construida para ese vacío;
    \item $D_t$: los documentos originales recuperados;
    \item $z_t$: la evidencia comprimida, que conserva solo lo relevante para el razonamiento actual.
\end{itemize}

\videofigure{figures/agentic-rag-mechanism.jpg}{El modelo usa marcadores de búsqueda especiales para pausar el reasoning, generar una query, invocar al recuperador e insertar de nuevo la evidencia en la cadena de razonamiento.}{00:54:42--00:56:22}

\videofigure{figures/reason-in-documents.jpg}{El módulo Reason-in-Documents analiza de forma independiente los documentos recuperados y extrae la información relevante, evitando que un documento largo interrumpa directamente la reasoning chain principal.}{00:56:22--00:58:48}

\begin{knowledgebox}{El compresor es, en esencia, un cuello de botella de información}
Debe conservar los hechos suficientes para sostener el siguiente paso del razonamiento y a la vez descartar los detalles irrelevantes. Comprimir de más pierde evidencia; comprimir de menos vuelve a traer ruido al modelo principal.
\end{knowledgebox}

\subsection{Worked example: un problema de química de varios pasos}

\videofigure{figures/search-o1-example.jpg}{En un problema de reacciones de varios pasos, Search-o1 primero localiza la estructura de un compuesto desconocido y después completa el conteo de átomos de carbono con la evidencia recuperada.}{00:58:48--01:00:34}

El vanilla reasoning simplemente adivina ante un compuesto desconocido; un agentic RAG común sí recupera documentos, pero un texto largo puede interrumpir el razonamiento original. Search-o1 divide el proceso en «detectar el vacío, consultar la estructura, depurar los hechos, continuar el cálculo», por lo que la respuesta final resulta más estable.

\subsection{Resumen de la sección}

Search-o1 convierte la recuperación de un procesamiento previo al inicio en una acción interna del reasoning, y añade una capa de compresión de documentos. Resuelve el problema de que las necesidades de conocimiento en una cadena de razonamiento larga se manifiestan de manera gradual.
